\section{Introduction}

Modern tourism often depends on multiple disconnected digital services. A traveler may need one application for maps, another for transit schedules, another for restaurant recommendations, and yet another source for local news or safety-related updates. As a result, users are forced to repeatedly switch contexts and manually combine information from several unrelated platforms. This fragmented experience is inconvenient, time-consuming, and especially difficult for tourists who are unfamiliar with the area they are visiting.

Urban Scout was designed to address this problem by providing a centralized system for tourism-related information. The project focuses on integrating key categories of data that are often used together in practice: places of interest, eateries, transit stations and schedules, localized news, safety alerts, incident reports, and itineraries. Rather than treating these pieces of information as separate resources, Urban Scout connects them through a relational database structure so that users can interact with them in a unified way.

The system created for this project is a mobile-friendly web application supported by a MySQL relational database, an Express/Node.js backend, and a React frontend. The primary end users are tourists, who can explore regions, browse places, check transit information, read local updates, report incidents, and organize their travel plans. A secondary user group consists of administrators, who are responsible for posting regional news and verifying incident reports. Through this design, Urban Scout demonstrates how database-driven integration can improve usability and reduce the burden of managing travel information across multiple disconnected applications.
